% Ortak önbölüm. pdflatex ve xelatex ile derlenir.
\usepackage{iftex}
\ifPDFTeX
  \usepackage[T1]{fontenc}
  \usepackage[utf8]{inputenc}
\else
  \usepackage{fontspec}
\fi

\usepackage[margin=1in]{geometry}
\usepackage{booktabs}
\usepackage{longtable}
\usepackage{array}
\usepackage{amsmath}
\usepackage{microtype}
\usepackage{graphicx}
\usepackage{xcolor}
\usepackage{caption}
\usepackage[hidelinks]{hyperref}
\usepackage{url}

\captionsetup{font=small,labelfont=bf}
\setlength{\tabcolsep}{5pt}
\renewcommand{\arraystretch}{1.12}

% Sayı sütunları için sağa dayalı, sabit genişlikli hücre.
\newcolumntype{R}[1]{>{\raggedleft\arraybackslash}p{#1}}
\newcolumntype{L}[1]{>{\raggedright\arraybackslash}p{#1}}

% Kod ve sütun adları.
\newcommand{\col}[1]{\texttt{\small #1}}

\hypersetup{
  pdfauthor={Serdar I. Caglar},
  pdfkeywords={speech corpus, text-to-speech, Turkish, forced alignment, data curation}
}
